\subsection{Hasil Pengujian}

Prototype awal CleanConnect yang ditampilkan pada Bab IV menjadi dasar untuk pelaksanaan evaluasi lanjutan. Pengujian dapat dilakukan dengan meminta calon pengguna mencoba alur utama aplikasi, seperti memilih layanan, menentukan lokasi, melihat estimasi biaya, memantau status petugas, mengatur jadwal rutin, dan memberikan rating setelah layanan selesai. Pada tahap draft ini, hasil pengujian formal belum dicantumkan karena belum tersedia data uji dari pengguna. Dokumentasi wawancara sebagai bukti pendukung pengumpulan kebutuhan pengguna disajikan pada lampiran.

\subsection{Kelebihan}

Berdasarkan rancangan solusi yang telah dijelaskan, CleanConnect memiliki beberapa kelebihan. Platform ini memudahkan pengguna menemukan jasa kebersihan rumah tangga tanpa bergantung sepenuhnya pada rekomendasi informal. Sistem rating dan review membantu pengguna menilai kredibilitas petugas berdasarkan pengalaman pengguna sebelumnya \cite{Pena-garcia2024,Id2018}. Selain itu, fitur estimasi biaya otomatis memberikan transparansi harga sebelum pemesanan dikonfirmasi, tracking petugas berbasis lokasi/GPS membantu pengguna mengetahui posisi dan estimasi waktu kedatangan petugas, sedangkan jadwal rutin memudahkan pengguna yang membutuhkan layanan kebersihan secara berkala \cite{Bernarto2022,Alvin2025}.

\subsection{Kekurangan}

CleanConnect juga memiliki beberapa keterbatasan yang perlu diperhatikan. Platform membutuhkan data petugas yang valid dan proses verifikasi yang baik agar kepercayaan pengguna dapat dibangun. Fitur tracking petugas bergantung pada akurasi layanan lokasi dan koneksi internet, sehingga kualitas pelacakan dapat menurun apabila jaringan tidak stabil \cite{Hawlitschek2016,Li2020}. Selain itu, estimasi biaya otomatis masih memerlukan parameter yang tepat dan rancangan sistem belum dapat dinilai secara menyeluruh sebelum dilakukan pengujian langsung kepada calon pengguna dan mitra petugas.
